\documentclass[11pt,a4paper]{article}
\usepackage{fontspec,amsmath,amssymb,amsthm,microtype}
\setmainfont{lmroman10-regular.otf}[BoldFont=lmroman10-bold.otf,ItalicFont=lmroman10-italic.otf,BoldItalicFont=lmroman10-bolditalic.otf]
\setmonofont{lmmono10-regular.otf}[Scale=MatchLowercase]
\usepackage[margin=25mm]{geometry}
\usepackage[unicode,hidelinks]{hyperref}
\setlength{\emergencystretch}{3em}
\setlength{\parskip}{0.4em}
\setlength{\parindent}{0pt}
\allowdisplaybreaks[2]
\raggedbottom
\title{\texorpdfstring{Construct spatial smoothness, finite Fisher information and finite actual entropy}{Construct spatial smoothness, finite Fisher information and finite actual entropy}}
\author{LangevinLean mathematical exposition}
\date{}
\begin{document}\maketitle
\noindent Source module: \texttt{Unadjusted\allowbreak{}Langevin/\allowbreak{}LSI/\allowbreak{}Gaussian\allowbreak{}Entropy.\allowbreak{}lean}. The accompanying HTML pairs this mathematical exposition with the exact Lean statements and records the scope of the formalization.


For \(t>0\), the actual frozen Gaussian update \(Y_t=X-t\nabla U(X)+\sqrt{2t}\,Z\) has the positive mixture density constructed in \texttt{Interpolation\allowbreak{}Density.\allowbreak{}lean}. This module proves its spatial \(C^1\) regularity by differentiating the displayed Gaussian integral. Bayes' posterior formula and Jensen's inequality then give the universal Lebesgue Fisher bound \(\int\|\nabla\log p_t\|^2p_t\le d/(2t)\). If the initial gradient has finite energy, the relative score \(\nabla\log r_t=\nabla U+\nabla\log p_t\) also has finite energy. The target LSI consequently proves finite actual KL at every positive time. No initial density or second moment of the state is assumed.


\section{\texorpdfstring{Bound the quadratic Gaussian envelope uniformly}{Bound the quadratic Gaussian envelope uniformly}}


For \(a>0\), put \(u=a\|z\|^2\). The scalar inequality \(u e^{-u}\le e^{-1}\le1\) gives

\[\|z\|^2e^{-a\|z\|^2}\le a^{-1}.\]

Lean multiplies the desired inequality by the positive number \(a\), rewrites its left side as \(u e^{-u}\), and applies \texttt{Real.\allowbreak{}mul\_\allowbreak{}exp\_\allowbreak{}neg\_\allowbreak{}le\_\allowbreak{}exp\_\allowbreak{}neg\_\allowbreak{}one}. The final cancellation uses \texttt{mul\_\allowbreak{}inv\_\allowbreak{}cancel\ensuremath{{}_{0}}} with the explicit nonzero witness from \texttt{ha}. This global bound later controls the conditional displacement variance even when the initial distribution has no moments.



{\raggedright\noindent Lean declarations: \texttt{norm\_\allowbreak{}sq\_\allowbreak{}mul\_\allowbreak{}exp\_\allowbreak{}neg\_\allowbreak{}norm\_\allowbreak{}sq\_\allowbreak{}le}.\par}

\section{\texorpdfstring{Bound the first Gaussian envelope without an initial moment}{Bound the first Gaussian envelope without an initial moment}}


The elementary estimate \(r\le1+r^2\), together with \(e^{-ar^2}\le1\) and the quadratic envelope bound, gives

\[\|z\|e^{-a\|z\|^2}\le1+a^{-1}.\]

The proof records the exponential bound as \texttt{hexp}, the quadratic bound as \texttt{hsq}, and \(r\le1+r^2\) as \texttt{hnorm}. Multiplication by the nonnegative exponential and addition then yield the displayed result. The bound is intentionally simple rather than optimal: a finite constant independent of the proposed state and mixture component is what dominated differentiation requires.



{\raggedright\noindent Lean declarations: \texttt{norm\_\allowbreak{}mul\_\allowbreak{}exp\_\allowbreak{}neg\_\allowbreak{}norm\_\allowbreak{}sq\_\allowbreak{}le}.\par}

\section{\texorpdfstring{Specify the real Gaussian kernel and its derivative map}{Specify the real Gaussian kernel and its derivative map}}


With precision \(a\), center \(c\), and the already normalized Gaussian constant \(N_a\), define

\[K_a(c,y)=N_a e^{-a\|y-c\|^2},\qquad D_yK_a(c,y)[v]=-2aK_a(c,y)\langle y-c,v\rangle.\]

\texttt{real\allowbreak{}Gaussian\allowbreak{}Kernel} is the ordinary real density. \texttt{real\allowbreak{}Gaussian\allowbreak{}Kernel\allowbreak{}Deriv} is its proposed Fréchet derivative as a continuous linear functional: the scalar \texttt{real\allowbreak{}Gaussian\allowbreak{}Kernel a c y * (-\allowbreak{}2 * a)} multiplies \texttt{inner\allowbreak{}SL \ensuremath{\mathbb R} (y -\allowbreak{} c)}. The precision parameter keeps this helper reusable; the ULA specialization is \(a=1/(4t)\). The next theorem proves that this proposed derivative is the actual derivative.



{\raggedright\noindent Lean declarations: \texttt{real\allowbreak{}Gaussian\allowbreak{}Kernel}, \texttt{real\allowbreak{}Gaussian\allowbreak{}Kernel\allowbreak{}Deriv}.\par}

\section{\texorpdfstring{Differentiate the kernel by the squared-\allowbreak{}norm and exponential chain rules}{Differentiate the kernel by the squared-norm and exponential chain rules}}


The map \(y\mapsto y-c\) has derivative the identity. The squared-norm derivative is \(v\mapsto2\langle y-c,v\rangle\). Multiplication by \(-a\), exponentiation, and multiplication by \(N_a\) therefore produce exactly

\[D_yK_a(c,y)[v]=-2aK_a(c,y)\langle y-c,v\rangle.\]

The local Lean proof \texttt{hd} is this sequence of \texttt{sub\_\allowbreak{}const}, \texttt{norm\_\allowbreak{}sq}, \texttt{const\_\allowbreak{}mul}, \texttt{exp}, and a final \texttt{const\_\allowbreak{}mul}. The final extensionality argument compares the two continuous linear maps at an arbitrary vector \texttt{v}; elementary scalar algebra proves they agree. No regularity of the potential is involved in this fixed-center calculation.



{\raggedright\noindent Lean declarations: \texttt{has\allowbreak{}FDeriv\allowbreak{}At\_\allowbreak{}real\allowbreak{}Gaussian\allowbreak{}Kernel}.\par}

\section{\texorpdfstring{Find a common integrable bound for all component derivatives}{Find a common integrable bound for all component derivatives}}


The norm of the inner-product functional \(v\mapsto\langle y-c,v\rangle\) equals \(\|y-c\|\). Hence

\[\|D_yK_a(c,y)\|=2aN_a\|y-c\|e^{-a\|y-c\|^2}\le2aN_a(1+a^{-1}).\]

Lean uses \texttt{norm\_\allowbreak{}smul} and \texttt{inner\allowbreak{}SL\_\allowbreak{}apply\_\allowbreak{}norm}, removes absolute values using positivity of \(N_a\) and the exponential, and applies the first Gaussian envelope bound. The resulting constant is independent of both \(y\) and \(c\). A probability mixture therefore integrates this bound without requiring any moment of its center map.



{\raggedright\noindent Lean declarations: \texttt{norm\_\allowbreak{}real\allowbreak{}Gaussian\allowbreak{}Kernel\allowbreak{}Deriv\_\allowbreak{}le}.\par}

\section{\texorpdfstring{Define the actual proposed mixture derivative}{Define the actual proposed mixture derivative}}


For \(c_t(x)=x-t\nabla U(x)\), define the continuous-linear-map-valued integral

\[P'_t(y)=\int D_yK_{1/(4t)}(c_t(x),y)\,d\mu(x).\]

This is \texttt{interpolation\allowbreak{}Lebesgue\allowbreak{}Density\allowbreak{}Deriv}. The symbol here denotes a spatial Fréchet derivative, not the time derivative of the density. The following domination lemmas establish its integrability, continuity, and equality with the actual spatial derivative of the explicit mixture. Keeping the definition as an integral of linear maps makes each directional evaluation correspond directly to the component kernel derivative.



{\raggedright\noindent Lean declarations: \texttt{interpolation\allowbreak{}Lebesgue\allowbreak{}Density\allowbreak{}Deriv}.\par}

\section{\texorpdfstring{Integrate each Gaussian component over the mixing probability law}{Integrate each Gaussian component over the mixing probability law}}


For a measurable center map \(c\), the function \(x\mapsto K_a(c(x),y)\) is measurable and satisfies \(0\le K_a(c(x),y)\le N_a\). Since \(\mu\) is a probability measure, the constant \(N_a\) is integrable, and so is the kernel.



Lean starts from \texttt{integrable\_\allowbreak{}const} and uses \texttt{Integrable.\allowbreak{}mono'}. Its first obligation is measurability, assembled from the center's measurability by \texttt{fun\_\allowbreak{}prop}. Its second obligation is the norm bound; positivity removes the absolute value, and \texttt{Real.\allowbreak{}exp\_\allowbreak{}le\_\allowbreak{}one\_\allowbreak{}iff} supplies the Gaussian peak bound.



{\raggedright\noindent Lean declarations: \texttt{integrable\_\allowbreak{}real\allowbreak{}Gaussian\allowbreak{}Kernel}.\par}

\section{\texorpdfstring{Integrate the continuous-\allowbreak{}linear-\allowbreak{}map-\allowbreak{}valued component derivative}{Integrate the continuous-linear-map-valued component derivative}}


The common derivative bound is a finite constant \(2aN_a(1+a^{-1})\). Therefore

\[\int\|D_yK_a(c(x),y)\|\,d\mu(x)<\infty.\]

The Lean proof again applies \texttt{Integrable.\allowbreak{}mono'} to an integrable constant, now with values in the Banach space of continuous linear functionals. Measurability follows from the explicit derivative formula; the second obligation is exactly \texttt{norm\_\allowbreak{}real\allowbreak{}Gaussian\allowbreak{}Kernel\allowbreak{}Deriv\_\allowbreak{}le}. This supplies the Bochner-integrability premise needed to identify the derivative of the mixture.



{\raggedright\noindent Lean declarations: \texttt{integrable\_\allowbreak{}real\allowbreak{}Gaussian\allowbreak{}Kernel\allowbreak{}Deriv}.\par}

\section{\texorpdfstring{Differentiate the mixture under the integral sign}{Differentiate the mixture under the integral sign}}


For each proposed state \(y\),

\[D_y\!\left(\int K_a(c(x),y)\,d\mu(x)\right)=\int D_yK_a(c(x),y)\,d\mu(x).\]

Lean applies \texttt{has\allowbreak{}FDeriv\allowbreak{}At\_\allowbreak{}integral\_\allowbreak{}of\_\allowbreak{}dominated\_\allowbreak{}of\_\allowbreak{}fderiv\_\allowbreak{}le} with the whole spatial domain as the common neighborhood. Kernel integrability supplies the value at the base point, derivative-kernel integrability supplies derivative measurability, the uniform Gaussian bound supplies domination, and \texttt{has\allowbreak{}FDeriv\allowbreak{}At\_\allowbreak{}real\allowbreak{}Gaussian\allowbreak{}Kernel} supplies pointwise differentiation. All these premises were proved from the actual kernel. In particular, a derivative of the mixture is not among the hypotheses.



{\raggedright\noindent Lean declarations: \texttt{has\allowbreak{}FDeriv\allowbreak{}At\_\allowbreak{}integral\_\allowbreak{}real\allowbreak{}Gaussian\allowbreak{}Kernel}.\par}

\section{\texorpdfstring{Prove continuity of the mixture derivative with the same domination}{Prove continuity of the mixture derivative with the same domination}}


For every fixed initial point \(x\), the derivative kernel \(y\mapsto D_yK_a(c(x),y)\) is continuous. Its operator norm is bounded by the same integrable constant for every \(y\). Dominated continuity therefore proves continuity of

\[y\longmapsto\int D_yK_a(c(x),y)\,d\mu(x).\]

The Lean theorem \texttt{continuous\_\allowbreak{}of\_\allowbreak{}dominated} requires measurability of every section, a common integrable bound, and almost-everywhere continuity of the component map. The preceding lemmas provide the first three facts, while the explicit formula and \texttt{fun\_\allowbreak{}prop} prove continuity. This continuity is what strengthens mere differentiability to \(C^1\).



{\raggedright\noindent Lean declarations: \texttt{continuous\_\allowbreak{}integral\_\allowbreak{}real\allowbreak{}Gaussian\allowbreak{}Kernel\allowbreak{}Deriv}.\par}

\section{\texorpdfstring{Combine the derivative and its continuity into a C1 theorem}{Combine the derivative and its continuity into a C1 theorem}}


A function is \(C^1\) exactly when it has a continuous Fréchet derivative. Apply this criterion to the derivative field \(y\mapsto\int D_yK_a(c(x),y)\,d\mu(x)\).



\texttt{cont\allowbreak{}Diff\_\allowbreak{}one\_\allowbreak{}iff\_\allowbreak{}has\allowbreak{}FDeriv\allowbreak{}At} reduces the Lean proof to the two preceding results: \texttt{continuous\_\allowbreak{}integral\_\allowbreak{}real\allowbreak{}Gaussian\allowbreak{}Kernel\allowbreak{}Deriv} gives continuity of this field, and \texttt{has\allowbreak{}FDeriv\allowbreak{}At\_\allowbreak{}integral\_\allowbreak{}real\allowbreak{}Gaussian\allowbreak{}Kernel} identifies it as the derivative at every point. The conclusion applies to every probability mixing measure, including singular ones.



{\raggedright\noindent Lean declarations: \texttt{cont\allowbreak{}Diff\_\allowbreak{}integral\_\allowbreak{}real\allowbreak{}Gaussian\allowbreak{}Kernel}.\par}

\section{\texorpdfstring{Apply Gaussian smoothing to the actual frozen-\allowbreak{}drift law}{Apply Gaussian smoothing to the actual frozen-drift law}}


The explicit real density of the update satisfies

\[p_t(y)=\int K_{1/(4t)}(x-t\nabla U(x),y)\,d\mu(x).\]

The proof first establishes this equality as \texttt{heq}, using the integral formula from \texttt{Interpolation\allowbreak{}Density.\allowbreak{}lean} and the scalar identity \(-\|z\|^2/(4t)=-(1/(4t))\|z\|^2\). The center map is measurable because the gradient is Lipschitz. The general Gaussian-mixture \(C^1\) theorem then applies with positive precision \(1/(4t)\). Thus spatial regularity is a consequence of the actual transition law, not an assumption about a chosen interpolation density.



{\raggedright\noindent Lean declarations: \texttt{interpolation\allowbreak{}Lebesgue\allowbreak{}Density\_\allowbreak{}cont\allowbreak{}Diff}.\par}

\section{\texorpdfstring{Identify the exact spatial derivative of the actual density}{Identify the exact spatial derivative of the actual density}}


The derivative field defined earlier is now certified:

\[D_yp_t(y)=P'_t(y)=\int D_yK_{1/(4t)}(x-t\nabla U(x),y)\,d\mu(x).\]

Lean applies the general dominated-differentiation theorem to the measurable frozen center and rewrites its mixture back to \texttt{interpolation\allowbreak{}Lebesgue\allowbreak{}Density}. The final \texttt{rfl} checks that the derivative integral is exactly \texttt{interpolation\allowbreak{}Lebesgue\allowbreak{}Density\allowbreak{}Deriv}. This identity is later used to define a jointly measurable score field while retaining its agreement with the actual spatial derivative.



{\raggedright\noindent Lean declarations: \texttt{has\allowbreak{}FDeriv\allowbreak{}At\_\allowbreak{}interpolation\allowbreak{}Lebesgue\allowbreak{}Density}.\par}

\section{\texorpdfstring{Recover C1 regularity of the potential from its Lipschitz gradient}{Recover C1 regularity of the potential from its Lipschitz gradient}}


The assumptions are differentiability of \(U\) and Lipschitz continuity of \(\nabla U\). The Riesz isometry identifies

\[DU(x)=\operatorname{Riesz}(\nabla U(x)).\]

Since the Riesz map and the gradient are continuous, \(x\mapsto DU(x)\) is continuous. The criterion \texttt{cont\allowbreak{}Diff\_\allowbreak{}one\_\allowbreak{}iff\_\allowbreak{}fderiv} now yields \(U\in C^1\). Lean rewrites the derivative field by \texttt{to\allowbreak{}Dual\_\allowbreak{}comp\_\allowbreak{}gradient} and composes the two continuous maps. No Hessian or twice-continuously-differentiable potential is assumed.



{\raggedright\noindent Lean declarations: \texttt{cont\allowbreak{}Diff\_\allowbreak{}potential\_\allowbreak{}of\_\allowbreak{}lipschitz\_\allowbreak{}gradient}.\par}

\section{\texorpdfstring{Establish spatial admissibility of the actual Gibbs-\allowbreak{}relative density}{Establish spatial admissibility of the actual Gibbs-relative density}}


The relative density is \(r_t(y)=Z e^{U(y)}p_t(y)\). The partition factor \(Z\) is constant in space, \(e^U\) is \(C^1\) because \(U\) is \(C^1\), and the actual Gaussian mixture \(p_t\) is \(C^1\). Therefore \(r_t\) is \(C^1\) at every positive time.



The Lean proof follows this product exactly using \texttt{cont\allowbreak{}Diff\_\allowbreak{}const.\allowbreak{}mul}, the potential theorem followed by \texttt{exp}, and \texttt{interpolation\allowbreak{}Lebesgue\allowbreak{}Density\_\allowbreak{}cont\allowbreak{}Diff}. Together with pointwise positivity and normalization from the density-construction module, this discharges the smooth test-class premises of \texttt{Is\allowbreak{}Log\allowbreak{}Sobolev}.



{\raggedright\noindent Lean declarations: \texttt{interpolation\allowbreak{}Relative\allowbreak{}Density\_\allowbreak{}cont\allowbreak{}Diff}.\par}

\section{\texorpdfstring{Apply the target LSI directly to the constructed interpolation law}{Apply the target LSI directly to the constructed interpolation law}}


Use the actual positive normalized \(C^1\) relative density as the LSI test. This gives

\[2\alpha\,\mathrm{KL}(\rho_t\mid\nu)\le I(\rho_t\mid\nu).\]

This theorem retains extended nonnegative values, so neither finiteness is silently assumed. Lean obtains the Gibbs probability instance from \texttt{h\allowbreak{}LSI.\allowbreak{}1}; \texttt{h\allowbreak{}LSI.\allowbreak{}2} is applied with the already proved positivity, \(C^1\) regularity, and relative probability normalization. Finally \texttt{tilted\allowbreak{}Measure\_\allowbreak{}interpolation\allowbreak{}Relative\allowbreak{}Density} replaces the induced tilted measure with the actual Gaussian interpolation law. The only functional-inequality hypothesis is the target's LSI.



{\raggedright\noindent Lean declarations: \texttt{interpolation\allowbreak{}Law\_\allowbreak{}log\allowbreak{}Sobolev}.\par}

\section{\texorpdfstring{Record the square-\allowbreak{}norm Jensen inequality for probability measures}{Record the square-norm Jensen inequality for probability measures}}


For a square-integrable random vector \(F\) under a probability measure, convexity of squared norm gives

\[\left\|\int F\,d\mu\right\|^2\le\int\|F\|^2\,d\mu.\]

Lean obtains convexity by applying \texttt{Convex\allowbreak{}On.\allowbreak{}pow} to convexity and nonnegativity of the norm, with exponent two. \texttt{map\_\allowbreak{}integral\_\allowbreak{}le} then applies Jensen on the whole vector space. \texttt{Mem\allowbreak{}Lp} supplies both integrability of \(F\) and of \(\|F\|^2\). This general lemma is used under the normalized Gaussian posterior at each proposed state.



{\raggedright\noindent Lean declarations: \texttt{norm\_\allowbreak{}integral\_\allowbreak{}sq\_\allowbreak{}le\_\allowbreak{}integral\_\allowbreak{}norm\_\allowbreak{}sq}.\par}

\section{\texorpdfstring{Define Bayes posterior on the initial mixture component}{Define Bayes posterior on the initial mixture component}}


At a fixed proposed point \(y\), let \(p(y)=\int K_a(c(x),y)\,d\mu(x)\). The posterior law of the initial component is

\[\lambda_y(dx)=\frac{K_a(c(x),y)}{p(y)}\,\mu(dx).\]

\texttt{gaussian\allowbreak{}Posterior} constructs this measure with \texttt{with\allowbreak{}Density}. The definition is total, while its intended probability interpretation uses \(a>0\), measurability of the center, and positivity of the mixing probability law. The following lemmas establish positivity of the denominator and normalization explicitly, before using Jensen.



{\raggedright\noindent Lean declarations: \texttt{gaussian\allowbreak{}Posterior}.\par}

\section{\texorpdfstring{Prove the real mixture denominator is everywhere positive}{Prove the real mixture denominator is everywhere positive}}


Every component \(K_a(c(x),y)\) is positive because its normalizer and exponential are positive. Its support as a function of \(x\) is therefore the whole initial-state space, which has probability one.



The Lean proof uses \texttt{integral\_\allowbreak{}pos\_\allowbreak{}iff\_\allowbreak{}support\_\allowbreak{}of\_\allowbreak{}nonneg}, with integrability from the common Gaussian peak bound. The local support equality \texttt{hs} converts the condition to positive mass of the whole space. Thus \(p(y)>0\) at every point, allowing Bayes normalization and the logarithmic score to be defined pointwise.



{\raggedright\noindent Lean declarations: \texttt{integral\_\allowbreak{}real\allowbreak{}Gaussian\allowbreak{}Kernel\_\allowbreak{}pos}.\par}

\section{\texorpdfstring{Check posterior normalization by cancellation of the actual mixture integral}{Check posterior normalization by cancellation of the actual mixture integral}}


The total posterior mass equals

\[\int\frac{K_a(c(x),y)}{p(y)}\,d\mu(x)=\frac{p(y)}{p(y)}=1.\]

Lean first records positive denominator \texttt{hp} and integrable kernel \texttt{hq}. The definition of \texttt{with\allowbreak{}Density} expresses mass as a nonnegative integral. \texttt{of\allowbreak{}Real\_\allowbreak{}integral\_\allowbreak{}eq\_\allowbreak{}lintegral\_\allowbreak{}of\allowbreak{}Real} converts it to the ordinary integral, justified by integrability and nonnegativity of the normalized kernel. \texttt{integral\_\allowbreak{}div} pulls out the denominator and \texttt{div\_\allowbreak{}self hp.\allowbreak{}ne'} performs the normalization.



{\raggedright\noindent Lean declarations: \texttt{gaussian\allowbreak{}Posterior\_\allowbreak{}is\allowbreak{}Probability}.\par}

\section{\texorpdfstring{Prove posterior displacement has finite conditional second moment}{Prove posterior displacement has finite conditional second moment}}


The posterior may arise from an initial law without moments. Nevertheless, at each fixed \(y\),

\[\int\|y-c(x)\|^2\,d\lambda_y(x)=\frac{1}{p(y)}\int K_a(c(x),y)\|y-c(x)\|^2\,d\mu(x)<\infty.\]

The quadratic Gaussian envelope bounds the numerator's integrand by \(N_a/a\). Lean's \texttt{integrable\_\allowbreak{}with\allowbreak{}Density\_\allowbreak{}iff\_\allowbreak{}integrable\_\allowbreak{}smul'} changes posterior integrability into the displayed weighted integral under \(\mu\). The proof \texttt{hsquare} dominates it by the integrable constant \(N_a/(ap(y))\). Measurability follows from the measurable center and absolute continuity of the posterior. The standard \(L^2\) characterization then supplies \texttt{Mem\allowbreak{}Lp} for the displacement.



{\raggedright\noindent Lean declarations: \texttt{gaussian\allowbreak{}Posterior\_\allowbreak{}mem\allowbreak{}Lp\_\allowbreak{}displacement}.\par}

\section{\texorpdfstring{Identify the mixture score as the posterior component-\allowbreak{}score mean}{Identify the mixture score as the posterior component-score mean}}


Dividing the mixture derivative by its positive value gives

\[\nabla\log p(y)=-2a\int(y-c(x))\,d\lambda_y(x).\]

Lean differentiates \(\log p\) using the actual mixture derivative and the nonzero denominator. The Riesz inverse commutes with the Bochner integral; \texttt{hgrad} identifies the inverse image of each component derivative with \(-2aK_a(c(x),y)(y-c(x))\). Expanding the posterior via \texttt{integral\_\allowbreak{}with\allowbreak{}Density\_\allowbreak{}eq\_\allowbreak{}integral\_\allowbreak{}to\allowbreak{}Real\_\allowbreak{}smul} gives the same integral divided by \(p(y)\). The final scalar cancellations establish the equality pointwise, rather than merely choosing an almost-everywhere score representative.



{\raggedright\noindent Lean declarations: \texttt{gradient\_\allowbreak{}log\_\allowbreak{}integral\_\allowbreak{}real\allowbreak{}Gaussian\allowbreak{}Kernel}.\par}

\section{\texorpdfstring{Use posterior Jensen to bound the pointwise mixture score energy}{Use posterior Jensen to bound the pointwise mixture score energy}}


Applying square-norm Jensen to the posterior score identity yields

\[p(y)\|\nabla\log p(y)\|^2\le4a^2\int K_a(c(x),y)\|y-c(x)\|^2\,d\mu(x).\]

The posterior is a probability measure, and its displacement is square integrable by the preceding lemmas; these are the exact hypotheses needed for Jensen. Lean records that inequality as \texttt{hj}, scales it by \(4a^2p(y)\), expands the posterior integral, and cancels \(p(y)\) using positivity. This gives a usable unnormalized density inequality whose right side can be integrated over the proposed state by Fubini.



{\raggedright\noindent Lean declarations: \texttt{gaussian\allowbreak{}Mixture\_\allowbreak{}score\_\allowbreak{}energy\_\allowbreak{}le}.\par}

\section{\texorpdfstring{Compute the exact Gaussian component displacement energy}{Compute the exact Gaussian component displacement energy}}


At precision \(1/(4t)\), the component density is the law of \(c+\sqrt{2t}Z\). Therefore

\[\int K_{1/(4t)}(c,y)\|y-c\|^2\,dy=\mathbb E\|\sqrt{2t}Z\|^2=2td.\]

The proof includes integrability, not only an integral identity. \texttt{gaussian\_\allowbreak{}proposal\_\allowbreak{}radial\_\allowbreak{}density} identifies the normalized density with the actual affine Gaussian law. Gaussian square integrability proves \texttt{hnoise}; \texttt{integrable\_\allowbreak{}map\_\allowbreak{}measure} transfers it to the proposed state, and the density-integrability equivalence gives the weighted Lebesgue integrand. The integral identity then follows from the map formula, \((\sqrt{2t})^2=2t\), and the already proved standard Gaussian moment \(\mathbb E\|Z\|^2=d\).



{\raggedright\noindent Lean declarations: \texttt{real\allowbreak{}Gaussian\allowbreak{}Kernel\_\allowbreak{}displacement\_\allowbreak{}moment}.\par}

\section{\texorpdfstring{Integrate posterior Jensen to obtain the universal smoothing Fisher bound}{Integrate posterior Jensen to obtain the universal smoothing Fisher bound}}


Integrate the pointwise score inequality, use Fubini, and insert the component displacement moment:

\[\int p(y)\|\nabla\log p(y)\|^2\,dy\le4\left(\frac1{4t}\right)^2\int2td\,d\mu=\frac d{2t}.\]

Lean first proves joint integrability of the density-weighted displacement energy as \texttt{hprod}, using \texttt{integrable\_\allowbreak{}prod\_\allowbreak{}iff} and the constant component moment. The right side of pointwise Jensen is consequently integrable. Continuity of the mixture score supplies measurability of the left integrand, so \texttt{Integrable.\allowbreak{}mono'} proves \texttt{h\allowbreak{}W}. Only after establishing both integrability statements does the proof apply \texttt{integral\_\allowbreak{}mono}, exchange the integrals, and evaluate the Gaussian moment. The bound is independent of the initial distribution's location, density, or state moments.



{\raggedright\noindent Lean declarations: \texttt{gaussian\allowbreak{}Mixture\_\allowbreak{}fisher\_\allowbreak{}bound}.\par}

\section{\texorpdfstring{Propagate finite gradient energy without a state second moment}{Propagate finite gradient energy without a state second moment}}


Write

\[\nabla U(Y_t)=\bigl(\nabla U(Y_t)-\nabla U(X)\bigr)+\nabla U(X).\]

The first term is in \(L^2\) because Lipschitzness bounds it by \(L\|Y_t-X\|\), and the actual frozen displacement has finite second moment whenever \(\nabla U(X)\in L^2\). The second term is in \(L^2\) by assumption. Lean uses \texttt{frozen\_\allowbreak{}gradient\_\allowbreak{}error\_\allowbreak{}energy} and the \(L^2\) characterization for the discrepancy, adds the pulled-back initial gradient, and transfers the sum to the output measure with \texttt{mem\allowbreak{}Lp\_\allowbreak{}map\_\allowbreak{}measure\_\allowbreak{}iff}. The initial position cancels from the displacement, which is why no \(\mathbb E\|X\|^2\) premise appears.



{\raggedright\noindent Lean declarations: \texttt{mem\allowbreak{}Lp\_\allowbreak{}gradient\_\allowbreak{}interpolation\allowbreak{}Law}.\par}

\section{\texorpdfstring{Transfer the smoothing Fisher estimate to the actual update law}{Transfer the smoothing Fisher estimate to the actual update law}}


The generic mixture Fisher theorem applies with \(c_t(x)=x-t\nabla U(x)\). Rewriting its integral through \(\rho_t=p_t\,dy\) gives

\[\nabla\log p_t\in L^2(\rho_t),\qquad\int\|\nabla\log p_t\|^2\,d\rho_t\le\frac d{2t}.\]

The local identity \texttt{heq} matches the actual density with the generic kernel integral. \texttt{hf} is the generic integrability-and-bound result. The density-law identity and \texttt{integrable\_\allowbreak{}with\allowbreak{}Density\_\allowbreak{}iff\_\allowbreak{}integrable\_\allowbreak{}smul'} convert its integrability to the actual law; the integral-with-density formula converts the numerical bound. Pointwise finiteness of the density ensures that real and extended density representations agree.



{\raggedright\noindent Lean declarations: \texttt{interpolation\allowbreak{}Lebesgue\_\allowbreak{}fisher\_\allowbreak{}bound}.\par}

\section{\texorpdfstring{Retain everywhere positivity for the real density}{Retain everywhere positivity for the real density}}


The extended mixture density is positive and finite at every point for \(t>0\). Its real conversion is consequently positive:

\[p_t(y)>0.\]

The proof is exactly \texttt{ENNReal.\allowbreak{}to\allowbreak{}Real\_\allowbreak{}pos}, with the two previously proved guards \texttt{interpolation\allowbreak{}Density\_\allowbreak{}pos} and \texttt{interpolation\allowbreak{}Density\_\allowbreak{}ne\_\allowbreak{}top}. This real positivity witness is used whenever Lean differentiates \(\log p_t\) or rewrites the logarithm of a product.



{\raggedright\noindent Lean declarations: \texttt{interpolation\allowbreak{}Lebesgue\allowbreak{}Density\_\allowbreak{}pos}.\par}

\section{\texorpdfstring{Split the relative score into potential and Lebesgue scores}{Split the relative score into potential and Lebesgue scores}}


Since \(r_t=Ze^Up_t\), positivity and finite normalization give

\[\log r_t=\log Z+U+\log p_t,\qquad\nabla\log r_t=\nabla U+\nabla\log p_t.\]

Lean first proves that the real partition function is positive and finite. The function equality \texttt{heq} follows from \texttt{Real.\allowbreak{}log\_\allowbreak{}mul} and \texttt{Real.\allowbreak{}log\_\allowbreak{}exp}. Spatial differentiability of \(p_t\) then justifies the logarithmic derivative; differentiability of \(U\) handles the other summand. The proof differentiates the displayed sum and uses linearity of the Riesz inverse. The normalization constant disappears from the gradient exactly because it is spatially constant.



{\raggedright\noindent Lean declarations: \texttt{gradient\_\allowbreak{}log\_\allowbreak{}interpolation\allowbreak{}Relative\allowbreak{}Density}.\par}

\section{\texorpdfstring{Construct finite relative-\allowbreak{}score energy from two proved L2 fields}{Construct finite relative-score energy from two proved L2 fields}}


The potential gradient remains in \(L^2(\rho_t)\), while the Gaussian smoothing bound gives \(\nabla\log p_t\in L^2(\rho_t)\). Their sum is therefore in \(L^2\), and the score decomposition identifies it with the actual relative score:

\[\nabla\log r_t\in L^2(\rho_t).\]

Lean obtains continuity of the Lebesgue score from \(C^1\) regularity of its positive density, converts the finite Fisher integral to \texttt{Mem\allowbreak{}Lp}, and adds it to \texttt{mem\allowbreak{}Lp\_\allowbreak{}gradient\_\allowbreak{}interpolation\allowbreak{}Law}. The final function extensionality uses \texttt{gradient\_\allowbreak{}log\_\allowbreak{}interpolation\allowbreak{}Relative\allowbreak{}Density} pointwise. This is the finite-energy fact used for the bounded logarithmic entropy tests.



{\raggedright\noindent Lean declarations: \texttt{interpolation\allowbreak{}Relative\allowbreak{}Score\_\allowbreak{}mem\allowbreak{}Lp}.\par}

\section{\texorpdfstring{Prove finite relative Fisher information before taking real values}{Prove finite relative Fisher information before taking real values}}


Square integrability of the actual relative score implies

\[I(\rho_t\mid\nu)=\int\|\nabla\log r_t\|^2\,d\rho_t<\infty.\]

Lean starts from the \(L^2\) score theorem and obtains the integrable squared norm. It unfolds \texttt{relative\allowbreak{}Fisher}, replaces the tilted measure by the actual interpolation law, and rewrites the nonnegative integral as \texttt{ENNReal.\allowbreak{}of\allowbreak{}Real} of the ordinary energy integral. Such an extended value is finite. This explicit proof prevents later uses of \texttt{to\allowbreak{}Real} from discarding an infinite Fisher information.



{\raggedright\noindent Lean declarations: \texttt{interpolation\allowbreak{}Relative\allowbreak{}Fisher\_\allowbreak{}ne\_\allowbreak{}top}.\par}

\section{\texorpdfstring{Derive finite actual KL from LSI and the constructed finite Fisher energy}{Derive finite actual KL from LSI and the constructed finite Fisher energy}}


At positive time, the constructed density is an admissible LSI test, so \(2\alpha\,\mathrm{KL}(\rho_t\mid\nu)\le I(\rho_t\mid\nu)\). With \(\alpha>0\) and finite Fisher information, the KL divergence must be finite.



Lean uses the target's LSI probability instance, \texttt{interpolation\allowbreak{}Relative\allowbreak{}Fisher\_\allowbreak{}ne\_\allowbreak{}top}, and \texttt{interpolation\allowbreak{}Law\_\allowbreak{}log\allowbreak{}Sobolev}. If KL were infinity, its product with the positive nonzero factor \texttt{ENNReal.\allowbreak{}of\allowbreak{}Real (2 * \ensuremath{\alpha})} would be infinity. The inequality would then force Fisher information to be infinity, contradicting the proved finiteness. No entropy continuity or initial density smoothness is assumed to obtain this positive-time finiteness.



{\raggedright\noindent Lean declarations: \texttt{interpolation\allowbreak{}Law\_\allowbreak{}kl\allowbreak{}Div\_\allowbreak{}ne\_\allowbreak{}top}.\par}

\end{document}
